\begin{table}[t!]
	\centering
	\footnotesize
	\caption{NIST SP\,800-22 for 8M ReRAM SHA-256 output ($p$-value;~\textbf{bold}~$=$ pass,~\textcolor{red}{red}~$=$~fail). (1)~Frequency, (2)~Block Freq., (3)~Runs, (4)~Longest Run, (5)~Binary Matrix Rank, (6)~Spectral, (7)~Non-Overlap.~Template, (8)~Overlap.~Template, (9)~Maurer\'s Universal, (10)~Linear Complexity, (11)~Serial, (12)~Approx.~Entropy, (13)~Cumulative Sums, (14)~Random Excursion, (15)~Random Excursion Variant.}
	\label{tab:nist22}
	\begin{tabular*}{\columnwidth}{@{\extracolsep{\fill}}crrrrrr}
		\toprule
		Test & \multicolumn{2}{c}{RT-A} & \multicolumn{2}{c}{RT-B} & \multicolumn{2}{c}{Hot} \\
		\cmidrule(lr){2-3}\cmidrule(lr){4-5}\cmidrule(lr){6-7}
		& R0 & R1 & R0 & R1 & R0 & R1 \\
		\midrule
		1 & \textbf{0.154} & \textbf{0.705} & \textbf{0.549} & \textbf{0.458} & \textbf{0.135} & \textbf{0.991} \\
		2 & \textbf{0.865} & \textbf{0.514} & \textbf{0.361} & \textbf{0.838} & \textbf{0.815} & \textbf{0.142} \\
		3 & \textbf{0.250} & \textbf{0.769} & \textbf{0.621} & \textbf{0.588} & \textbf{0.334} & \textbf{0.907} \\
		4 & \textbf{0.825} & \textbf{0.872} & \textbf{0.515} & \textbf{0.458} & \textbf{0.927} & \textbf{0.193} \\
		5 & \textbf{0.574} & \textbf{0.847} & \textbf{0.487} & \textbf{0.797} & \textbf{0.141} & \textbf{0.010} \\
		6 & \textbf{0.136} & \textbf{0.503} & \textbf{0.609} & \textbf{0.548} & \textbf{0.489} & \textbf{0.801} \\
		7 & \textbf{1.000} & \textbf{1.000} & \textbf{1.000} & \textbf{1.000} & \textbf{1.000} & \textbf{1.000} \\
		8 & \textbf{0.796} & \textbf{0.378} & \textbf{0.059} & \textbf{0.600} & \textbf{0.878} & \textbf{0.758} \\
		9 & \textbf{0.118} & \textbf{0.421} & \textbf{0.488} & \textbf{0.112} & \textbf{0.480} & \textbf{0.076} \\
		10 & \textbf{0.767} & \textbf{0.664} & \textbf{0.210} & \textbf{0.223} & \textbf{0.510} & \textbf{0.803} \\
		11 & \textbf{0.123} & \textbf{0.707} & \textbf{0.141} & \textbf{0.074} & \textbf{0.036} & \textbf{0.460} \\
		12 & \textbf{0.123} & \textbf{0.754} & \textbf{0.319} & \textbf{0.075} & \textbf{0.036} & \textbf{0.835} \\
		13 & \textbf{0.201} & \textbf{0.820} & \textbf{0.620} & \textbf{0.500} & \textbf{0.134} & \textbf{0.605} \\
		14 & \textbf{0.087} & \textbf{0.099} & \textbf{0.257} & \textbf{0.142} & \textbf{0.036} & \textbf{0.126} \\
		15 & \textbf{0.134} & \textbf{0.279} & \textbf{0.179} & \textbf{0.177} & \textbf{0.315} & \textbf{0.211} \\
		\midrule
		\textbf{Pass} & 15/15 & 15/15 & 15/15 & 15/15 & 15/15 & 15/15 \\
		\bottomrule
	\end{tabular*}
\end{table}
\section{Results}
\subsection{NIST SP800-90B Entropy Estimates}
Table~\ref{tab:nist90b} reports non-IID min-entropy estimates from $1\mathrm{M}$ samples per source across the nine NIST-90B estimators. We now evaluate the reultes for each entropy source independently.

\paragraph{4M Entropy Sources}
All three 4Mb conditions produce $\hat{H}_{\min} < 0.01$bits/symbol (0.0060, 0.0051, and 0.0079 for RT-A, RT-B, and 55\textdegree{}C, respectively).  The compression estimator (Test~3) returns the lowest values (0.0004 and 0.0046 for RT-A/B), confirming that latency sequences are highly compressible.  Moreover, this failure is not isolated to a single estimator: \emph{all nine} SP800-90B estimators return near-zero values for every 4Mb condition, ruling out any estimator-specific artifact and confirming a fundamental collapse of source entropy under the runtime measurement regime.

The near-zero estimates are consistent with an underlying latency distribution that has collapsed to a near-deterministic, narrow support.  The compression estimator value of $4\times10^{-4}$bits/symbol for RT-A is especially diagnostic: this approaches the theoretical minimum for a length-$1\mathrm{M}$ sequence over a 256-bin alphabet, and is consistent with the sequence being almost entirely concentrated in a single latency bin across all measurement rounds.  The Most Common Value estimator (Test~1) corroborates this, reporting 0.0074bits/symbol,a value dictated by the $\log_2$ of the dominant bin occupancy rather than any meaningful distributional spread.

The elevated-temperature condition (55\textdegree{}C) yields a marginal improvement to $\hat{H}_{\min} = 0.0079$bits/symbol.  While thermal activation can increase filament stochasticity by reducing the energy barrier for vacancy migration~\cite{yu2012switching}, this effect is insufficient to rescue the 4M entropy source: even at 55\textdegree{}C the result remains three orders of magnitude below what is required to seed a practical conditioner. The temperature gain is also notably smaller on the 4M than on the 8M device (compare the Hot column in Table~\ref{tab:nist90b}), suggesting the dominant limitation is not thermal but structural.

We attribute the runtime entropy collapse to structural bursting effects. In the characterization performed in Section \ref{sec:opt-burst} we explored all bursts that maintained word-alignment. During characterization, the full sweep of burst word-aligned burst offsets explored the bursting space constrained only by burst continuity. Under the runtime sequence, performing fixed-width bursts at burst-alisgned offsets within the reserved 256-byte region, nearly all write latency variance was removed.

\paragraph{8M Entropy Sources}
This interpretation is strengthened by the qualitative contrast with the 8Mb device, which is subjected to the same fixed-sequence methodology yet retains substantial entropy. That is the characterization The 8Mb crossbar is driven with 32-word bursts of $D_{\min}$-weight codewords, whereas the 4Mb operates on 4-word bursts.  The longer burst aggregates more switching events per transaction, increasing the thermal heterogeneity experienced across bitcells within a single burst and reducing the degree to which any individual filament's history can dominate the aggregate latency measurement.  This suggests a minimum burst complexity below which the single-stimulus feedback loop is not broken by the accumulated stochastic contribution of other cells, and the 4M 4-word burst falls below that threshold.

The 8Mb device produces $\hat{H}_{\min}$ of 2.4973 (RT-A), 1.6500 (RT-B), and 1.9844 (55\textdegree{}C)  bits/symbol.  The $t$-Tuple estimator (Test~4) is binding across all three conditions (2.8370, 1.7358, 2.2754), consistent with short-range thermal coupling between adjacent writes.  The 0.34bit restart penalty on RT-A reflects modest startup bias as the filament distribution converges after reset. Device-to-device variation (0.85bits/symbol between RT-A and RT-B) underscores the necessity of per-chip characterization.  The 55\textdegree{}C intermediate result (1.9844bits/symbol) suggests thermal activation provides a real but bounded boost to filament stochasticity.

\begin{table}[t]
	\centering

	\footnotesize \caption{NIST SP\,800-90B Non-IID entropy estimates.
		(1)~Most Common Value, (2)~Collision, (3)~Compression, (4)~$t$-Tuple,
		(5)~LRS, (6)~Multi-MCV, (7)~Lag Prediction, (8)~MultiMMC, (9)~LZ78Y.}
	\label{tab:nist90b} \setlength{\tabcolsep}{10pt}
	\begin{tabular*}{\columnwidth}{@{\extracolsep{\fill}}crrrrrr}
		\toprule
		Test & \multicolumn{3}{c}{4\,Mb} & \multicolumn{3}{c}{8\,Mb} \\
		\cmidrule(lr){2-4}\cmidrule(lr){5-7}
		& RT-A & RT-B & Hot & RT-A & RT-B & Hot \\
		\midrule
		1 & 0.0074 & 0.0074 & 0.1426 & 2.9477 & 1.8743 & 2.3675 \\
		2 & 0.0074 & 0.0075 & 0.0812 & 1.0000 & 1.0000 & 1.0000 \\
		3 & 0.0004 & 0.0046 & 0.0572 & 0.5618 & 0.7271 & 0.3192 \\
		4 & 0.0060 & 0.0051 & 0.0079 & 2.8370 & 1.7358 & 2.2754 \\
		5 & 0.0142 & 0.0116 & 0.0176 & 3.7061 & 2.5852 & 3.1420 \\
		6 & 0.0074 & 0.0074 & 0.0116 & 2.9550 & 1.8786 & 2.3749 \\
		7 & 0.0094 & 0.0077 & 0.0116 & 3.6370 & 2.3875 & 3.0965 \\
		8 & 0.0074 & 0.0075 & 0.0116 & 2.9596 & 1.8753 & 2.3702 \\
		9 & 0.0074 & 0.0075 & 0.0116 & 2.9626 & 1.8769 & 2.3745 \\
		\midrule
		$\min$ & \textbf{0.0060} & \textbf{0.0051} & \textbf{0.0079} & \textbf{2.8370} & \textbf{1.7358} & \textbf{2.2754} \\
		\midrule
		$\hat{H}_{\min}$ & \textbf{0.0060} & \textbf{0.0051} & \textbf{0.0079} & \textbf{2.4973} & \textbf{1.6500} & \textbf{1.9844} \\
		\bottomrule
	\end{tabular*}
\end{table}

\subsection{Relationship to Prior Work}
It is important to note that the 4M failure is not a disqualification of COTS ReRAM entropy harvesting in general, nor does it imply that the device contains no stochastic dynamics. Arul et al.~\cite{tolga2024trng} obtained non-trivial min-entropy from the same Fujitsu/RAMXEED device family by applying randomized write patterns followed by aggressive Von Neumann debiasing. We note that the structural latency component removed in Section \ref{sec:remove-structural} was on the order of tenths of a millisecond, while the timer LSB used by authors of \cite{tolga2024trng} likely operated on the order of MHz. The present work confirms that under the more conservative fixed-sequence regime demanded by a stationary entropy source evaluation, the 4M device cannot sustain an exploitable entropy rate. The 4Mb result points to a key observation that characterization-phase entropy measurements conducted over diverse stimuli can substantially overestimate the operational entropy of a ReRAM-based TRNG where the write pattern itself must be considered as a separate operating point.

\subsection{NIST SP800-22 and PractRand}
Given the extremely low entropy of the 4M under a small wear-leveled burst sequence, we exclude it from our TRNG analysis. While Figure \ref{fig:eval-entropy} contains a limited dataset and subsequent symbol and output encoding from the NIST-90B dataset, capturing sufficient samples to perform the following tests on the 4M would be intractable given the >1ms burst write time of the 4M.

Table~\ref{tab:nist22} reports SP800-22 results for 256KB of SHA-256-conditioned output cap from all three 8Mb sources across two independent runs each.  All 15 tests pass for every source-run combination (15/15 across all six columns). $p$-values are well-distributed across $(0,1)$; the lowest observed is 0.01 (Binary Matrix Rank, 55\textdegree{}C-R1), which nearly meets the the $\alpha = 0.01$ rejection threshold. All three streams additionally pass PractRand through 256KB with no anomalies at any evaluation point, confirming that no artifacts from the repeating $S_{\mathrm{opt}}$ sequence or wear-leveling pattern persist through conditioning. Together, these results validate the correct implementation of the conditioning pipeline; they do not independently attest to raw source entropy, which is established by SP800-90B above. The final TRNG achieves a throughput of 1.7kbits/sec. 

\section{Conclusion}
We have presented an ECC-aware framework for entropy harvesting from COTS ReRAM memories and demonstrated its application to two commercial device families.  By leveraging reverse-engineered ECC structure to isolate codeword-level switching compositions, we suppress the deterministic timing perturbations that corrupt prior TRNG claims and construct wear-leveled burst sequences whose min-entropy can be rigorously evaluated under NIST SP800-90B requirements. We demonstrate a complete pipeline including a burn-in phase to improve the baseline entropy of the dedicated entropy source region.
The 8M RAMXEED device achieves $\hat{H}_{\min}$ between 1.65 and 2.50bits/symbol across all tested conditions, passes all 15 NIST SP800-22 tests and PractRand at 256KB, and delivers SHA-256-conditioned output at 1.7kbit/s, a throughput gain of 17x over prior Von Neumann-based approaches. The 4Mb device fails to sustain entropy under the same fixed-sequence regime despite exhibiting measurable variability during characterization, exposing a stimulus-conditioned stochasticity effect that invalidates characterization-only evaluations. We demonstrate that ECC knowledge is a necessary but not sufficient condition for an efficient COTS ReRAM TRNG. The underlying crossbar architecture must also not influence the timing characteristics. Runtime evaluation under the intended operational sequence is required as entropy estimates derived from diverse characterization stimuli may not be representative of deployment performance.